\scanpage{440}
\mspage{352}
D'ailleurs, $\underline T$ étant str. pleine dans $N$, $\underline S$ est str. pleine dans $M$, et il reste à vérifier que le foncteur précédent est ess. surjectif. Soit donc $x\in\mathrm{Ob}\,\underline S$. On aura, grâce à b)
\[
x=\underset{I}{\underrightarrow{\lim}}\,x_i
\qquad\left|\begin{aligned}
 &x_i\in\deleted{\unsure{S}}M_{\deleted{\unsure{<}}\pi}\\
 &I=M_\pi/x
\end{aligned}\right. ,
\]
\deleted{et \unsure{de} foncteur $i\longmapsto x_i$, $I\longrightarrow M_\pi$}

D'ailleurs, le système inductif précédent est $\pi_0$-$\pi$-adapté, i.e (p. 339)

1) $I$ équiv. à une petite catégorie, $I$ grde devant $\pi$ (résulte du fait que $\pi$ est utile pour $M$, cf \textcircled{b}.

2) les limites inductives grdes devant $\pi_0$ et de cardinal $\leq\unsure{\pi|2}$ existent dans $I$ \inserted{résulte de $L_{\pi_0}$ sur $M$ [cf. \textcircled{d}]}

3) le foncteur $\xi:I\unsure{\text{---}|\longrightarrow}M$ y commute.

4) Les $x_i$ sont dans $M_\pi$.

Comme $f$ commute aux $\underrightarrow{\lim}$ grdes devant $\pi_0$, \deleted{\unsure{et}} (cf \textcircled{d}), a fortiori à celles grdes devant $\pi$, on \unsure{a|aura} donc
\[
f(x)\simeq\underset I{\underrightarrow{\lim}}\,f(x_i),
\]
je dis que le syst. ind. $(f(x_i))_{i\in I}$ de $N$ est $\pi_0$-$\pi$-adapté. Les conditions 1), 2) \unsure{sont en évidence} (ne concernant que $I$), le foncteur $I\longrightarrow N$ est le composé
% slanted
\margin{Cette commutation devrait avoir un nom, \unsure{disons|la} condition $L(\pi_0,\pi)$ sur un foncteur}
\diagram{\begin{tikzcd}[ampersand replacement=\&,column sep=small] I\arrow[r,"\xi"]\&M\&\overset{f}{\unsure{\text{---}|\longrightarrow}}\&N \end{tikzcd}}, donc commute aux $\underrightarrow{\lim}$ grdes devant $\pi_0$ et de cardinal $\leq\pi$, puisqu'il en est ainsi de $\xi$ et de $f$, enfin les $f(x_i)$ sont ds $N_\pi$ puisque par \textcircled{c} $f(M_{\unsure{\pi|\pi_0}})\subset N_\pi$.
% archive sheet number: 439
